% \iffalse meta-comment
%
% hit-caption.dtx 是各个类共用的双语题注
%
% \fi
%
% \section{共享双语题注}
%
% \changes{v3.2a}{0000/00/00}{双语题注抽出来给各个类共用，原先每个类各写一遍}
% \changes{v3.2a}{0000/00/00}{英文题注改成 \cs{caption} 的尾随可选参数，
%   \cs{bicaption} 降为兼容写法}
% 中英文对照的题注原先要换一个命令 \cs{bicaption}，用户得记住「这张图要不要双语」
% 并挑命令。现在写法只有一个：
% \begin{latex}
% \caption{中文题}
% \caption{中文题}[English caption]
% \caption[目录用短标题]{中文题}[English caption]
% \end{latex}
% 给了尾随的 \oarg{英文题} 就排双语，不给就单语。目录用的短标题仍是 \cs{caption}
% 原本那个前置可选参数，两者位置不同，不会混。
%
% \option{style/bilingual-caption} 是一键关掉所有双语题注的开关，默认开。关掉之后
% 英文题一概不排，同一份稿子投不同用途时不必逐条删。没做成按学位给默认值：写了就出、
% 不写就不出已经够用，多一层默认反而要用户去猜自己这一档是开是关。
%
% 键自己声明，没走 \cs{\_\_hit\_style\_pair:nn}：那个在 \file{hit-meta.dtx} 里，
% 而本模块在 \file{hithesis.ins} 的名单里排在它前面。
%
% \cs{bicaption} 仍然可用：\pkg{ccaption} 的五参数原形与本模板之前的三、四参数写法
% 都接得住，只是各提示一次，v4.1 移除。
%
%    \begin{macrocode}
%<*shared-caption>
%    \end{macrocode}
%
% \begin{macro}{\hit@setup@caption}
% 接管 \cs{caption} 要等到 \texttt{begindocument/end}。\pkg{caption}（由
% \pkg{subcaption} 带进来）和 \pkg{ccaption} 都在 \cs{AtBeginDocument} 里重定义
% \cs{caption}，在类文件里接管会被它们覆盖掉——踩过一次，接管看似成功，
% \cs{meaning}\cs{caption} 到了正文里已经是 \pkg{caption} 的定义了。
% \texttt{begindocument/end} 排在所有 \cs{AtBeginDocument} 之后。
%
% \cs{\_\_hit\_bicaption:w} 存的是 \pkg{ccaption} 的原形，这个在类文件里存就行；
% \cs{hit@original@caption} 存的是接管前那条完整的 \cs{caption} 链。
%    \begin{macrocode}
\bool_new:N \g__hit_bilingual_caption_bool
\bool_gset_true:N \g__hit_bilingual_caption_bool
\keys_define:nn { hit / style }
  {
    bilingual-caption .bool_gset:N = \g__hit_bilingual_caption_bool ,
    bilingual-caption .default:n = { true } ,
  }
\msg_new:nnn { hithesis } { bicaption-deprecated }
  {
    \iow_char:N \\ bicaption~已不推荐，英文题注改成~\iow_char:N \\ caption~的
    尾随可选参数： \\
    \iow_char:N \\ caption{中文题}[English~caption]。 \\
    旧写法将在~v4.1（2027~年主要版本）中移除。
  }
\msg_new:nnn { hithesis } { bicaption-five-args }
  {
    \iow_char:N \\ bicaption~的五参数写法已不推荐。 \\
    英文图表名请改用~\iow_char:N \\ hitsetup{const={figure-prefix-en=...}}~设置，
    题注写成~\iow_char:N \\ caption{中文题}[English~caption]。 \\
    旧写法将在~v4.1（2027~年主要版本）中移除。
  }
%    \end{macrocode}
%
% 排双语题注：走 \pkg{ccaption} 的原形，英文图表名取 \option{const} 族里的
% \option{figure-prefix-en} 与 \option{table-prefix-en}，用户不必每次手写
% \texttt{Fig.} 加负细空格去抵 \cs{fnum@figure} 插的那个空格。
%    \begin{macrocode}
\cs_new_protected:Npn \hit@bilingual@caption #1#2#3
  {
    \tl_if_novalue:nTF {#1}
      { \__hit_bicaption:w {} {#2} { \use:c { hit@ \@captype @prefix@en } } {#3} }
      { \__hit_bicaption:w [#1] {} {#2} { \use:c { hit@ \@captype @prefix@en } } {#3} }
  }
\cs_new_protected:Npn \hit@single@caption #1#2
  {
    \tl_if_novalue:nTF {#1}
      { \hit@original@caption {#2} }
      { \hit@original@caption [#1] {#2} }
  }
%    \end{macrocode}
%
% \cs{caption} 的接管。给了尾随的英文题、并且开关没关死，才走双语。
%    \begin{macrocode}
\cs_new_protected:Npn \hit@setup@caption
  {
    \cs_set_eq:NN \__hit_bicaption:w \bicaption
    \AddToHook { begindocument/end }
      {
    \cs_set_eq:NN \hit@original@caption \caption
    \RenewDocumentCommand \caption { o +m o }
      {
        \IfNoValueTF {##3}
          { \hit@single@caption {##1} {##2} }
          {
            \bool_if:NTF \g__hit_bilingual_caption_bool
              { \hit@bilingual@caption {##1} {##2} {##3} }
              { \hit@single@caption {##1} {##2} }
          }
      }
%    \end{macrocode}
%
% \cs{bicaption} 的兼容层。三种写法靠「后面还跟不跟一个花括号组」区分：跟着的是
% \pkg{ccaption} 的五参数原形，抓走那一组照原形排；不跟就是本模板之前的三、四参数
% 写法。探之前先跳空格，与 \cs{@ifnextchar} 的惯例一致——旧示例里那五个参数是连着
% 写的，但用户可能换了行。题注后面正常不会紧跟花括号组。
%    \begin{macrocode}
    \RenewDocumentCommand \bicaption { o +m +m +g }
      {
        \peek_meaning_ignore_spaces:NTF \c_group_begin_token
          { \__hit_bicaption_legacy:nnnn {##1} {##2} {##3} {##4} }
          {
            \msg_warning:nn { hithesis } { bicaption-deprecated }
            \IfNoValueTF {##4}
              { \hit@bilingual@caption {##1} {##2} {##3} }
              { \hit@bilingual@caption {##1} {##3} {##4} }
          }
      }
      }
  }
\cs_new_protected:Npn \__hit_bicaption_legacy:nnnn #1#2#3#4#5
  {
    \msg_warning:nn { hithesis } { bicaption-five-args }
    \tl_if_novalue:nTF {#1}
      { \__hit_bicaption:w {#2} {#3} {#4} {#5} }
      { \__hit_bicaption:w [#1] {#2} {#3} {#4} {#5} }
  }
%    \end{macrocode}
% \end{macro}
%
%    \begin{macrocode}
%</shared-caption>
%    \end{macrocode}
%
% \Finale
